% Fluxograma do processo de pesquisa psicométrica
% Requer: tikz, positioning, arrows.meta, shapes.geometric, calc, fit, backgrounds
% (já carregados no preâmbulo da tese)

\definecolor{violetA}{HTML}{6D28D9}
\definecolor{violetB}{HTML}{8B5CF6}
\definecolor{roseA}{HTML}{BE185D}
\definecolor{emeraldA}{HTML}{047857}
\definecolor{slateG}{HTML}{475569}
\definecolor{amberA}{HTML}{D97706}
\definecolor{amberC}{HTML}{FEF3C7}

\begin{tikzpicture}[
  phase/.style={
    fill=#1!12, draw=#1!60,
    font=\sffamily\bfseries\small,
    minimum width=2.0cm, minimum height=0.8cm,
    align=center, rounded corners=3pt,
    line width=0.6pt
  },
  stepbox/.style={
    draw=#1!50, fill=#1!6,
    font=\small, align=center,
    minimum width=6.0cm, minimum height=0.9cm,
    rounded corners=3pt, text width=5.8cm,
    line width=0.6pt
  },
  metricbox/.style={
    draw=amberA!60, fill=amberC!40,
    font=\small, align=center,
    minimum width=4.6cm, minimum height=0.9cm,
    rounded corners=3pt, text width=4.4cm,
    line width=0.6pt, dashed
  },
  outputbox/.style={
    draw=#1!70, fill=#1!15,
    font=\small\bfseries, align=center,
    minimum width=6.0cm, minimum height=1.0cm,
    rounded corners=4pt, text width=5.8cm,
    line width=1pt
  },
  arr/.style={-{Stealth[length=5pt]}, thick, slateG!70},
  darr/.style={-{Stealth[length=4pt]}, thin, slateG!40, dashed},
  node distance=0.50cm
]

%% ═══ ETAPA 1: População e Amostra ═══
\node[phase=slateG] (ph1) {ETAPA~1\\[-2pt]{\scriptsize Amostra}};

\node[stepbox=slateG, right=0.6cm of ph1] (pop)
  {Defini\c{c}\~ao da popula\c{c}\~ao e amostra\\[-2pt]
   {\scriptsize Territ\'orio Semi\'arido NE~II $\cdot$ Comunidades quilombolas (Jeremoabo)}\\[-2pt]
   {\scriptsize $n = 30$--$40$ (Linacre, 1994) $\cdot$ Crit\'erios de inclus\~ao}};

%% ═══ ETAPA 2: Construção do Instrumento ═══
\node[phase=violetA, below=0.8cm of ph1] (ph2)
  {ETAPA~2\\[-2pt]{\scriptsize Instrumento}};

\node[stepbox=violetB, right=0.6cm of ph2] (instr)
  {Constru\c{c}\~ao do instrumento de pesquisa\\[-2pt]
   {\scriptsize 8 dimens\~oes ($V_1$--$V_8$) $\cdot$ Likert 5 pontos + perguntas abertas}\\[-2pt]
   {\scriptsize 35--45 itens $\cdot$ Derivado do WOCAT-SLM-QBR (Cap.~3)}};
\draw[arr] (pop) -- (instr);

%% ═══ ETAPA 3: Validação de Conteúdo ═══
\node[phase=roseA, below=0.8cm of ph2] (ph3)
  {ETAPA~3\\[-2pt]{\scriptsize Valida\c{c}\~ao}};

\node[stepbox=roseA, right=0.6cm of ph3] (valid)
  {Valida\c{c}\~ao de conte\'udo\\[-2pt]
   {\scriptsize Painel de ju\'{\i}zes $\cdot$ CVC $\geq$ 0,80}\\[-2pt]
   {\scriptsize Kappa de Fleiss $\cdot$ Avalia\c{c}\~ao \`as cegas}};
\draw[arr] (instr) -- (valid);

\node[metricbox, right=1.4cm of valid] (metr3)
  {Clareza de linguagem\\[-2pt]
   Pertin\^encia por dimens\~ao\\[-2pt]
   Relev\^ancia discriminante};
\draw[darr] (valid) -- (metr3);

%% ═══ ETAPA 4: Coleta de Dados ═══
\node[phase=emeraldA, below=0.8cm of ph3] (ph4)
  {ETAPA~4\\[-2pt]{\scriptsize Coleta}};

\node[stepbox=emeraldA, right=0.6cm of ph4] (coleta)
  {Coleta de dados em campo\\[-2pt]
   {\scriptsize 2--3 entrevistadores treinados $\cdot$ Manual de orienta\c{c}\~ao}\\[-2pt]
   {\scriptsize Supervis\~ao sistem\'atica $\cdot$ Relat\'orios padronizados}};
\draw[arr] (valid) -- (coleta);

%% ═══ ETAPA 5: Validação de Construto (Cap. 3) ═══
\node[phase=violetA, below=0.8cm of ph4] (ph5)
  {ETAPA~5\\[-2pt]{\scriptsize AFC}};

\node[stepbox=violetB, right=0.6cm of ph5] (afc)
  {Valida\c{c}\~ao de construto (Cap\'{\i}tulo~3)\\[-2pt]
   {\scriptsize AFC: CFI $\geq$ 0,90, RMSEA $\leq$ 0,08}\\[-2pt]
   {\scriptsize Confiabilidade: $\omega$ de McDonald $\geq$ 0,70}};
\draw[arr] (coleta) -- (afc);

\node[metricbox, right=1.4cm of afc] (metr5)
  {Efeitos de piso e teto\\[-2pt]
   Padr\~ao de respostas\\[-2pt]
   Modelo tripartite};
\draw[darr] (afc) -- (metr5);

%% ═══ ETAPA 6: Calibração TRI (Cap. 4) ═══
\node[phase=roseA, below=0.8cm of ph5] (ph6)
  {ETAPA~6\\[-2pt]{\scriptsize TRI}};

\node[stepbox=roseA, right=0.6cm of ph6] (tri)
  {Calibra\c{c}\~ao TRI (Cap\'{\i}tulo~4)\\[-2pt]
   {\scriptsize Modelo de Cr\'edito Parcial (Masters)}\\[-2pt]
   {\scriptsize Escores latentes $\theta_n$ (escala Rasch) $\cdot$ Winsteps 5.5.2}};
\draw[arr] (afc) -- (tri);

%% ═══ SAÍDA: ISB ═══
\node[outputbox=emeraldA, below=0.8cm of tri] (isb)
  {\'Indice de Salvaguarda Biocultural (ISB)\\[-2pt]
   {\scriptsize Infer\^encia fuzzy Mamdani $\cdot$ Escala 0--100}\\[-2pt]
   {\scriptsize Dossi\^es para inst\^ancias regulat\'orias}};
\draw[arr] (tri) -- (isb);

\end{tikzpicture}
